\section{Generalities}\label{sec:generalities}

\begin{itemize}
\item We write \(\R^{d + 1}_+ \eqd \R^d \times (0, \infty)\).
\item We write \(\langle x \rangle \eqd \sqrt{1 + | x |^2}\).
\item All \(L^p (\R^{d + 1}_+)\)-norms below are taken with respect to \(\dd x \frac{\dd t}{t}\).
\item We write \(\{\alpha\} \eqd \alpha - \lfloor \alpha \rfloor \in [0, 1)\).
\item A set \(\mathcal{B} \subset \Sch (\R^d)\) is bounded in \(\Sch (\R^d)\) if for every nonnegative integer \(M\),
\[
\sup_{\varphi \in \mathcal{B}} \sup_{z \in \R^d}
\sup_{| \beta | \leq M}
\langle z \rangle^M | \partial^{\beta} \varphi (z) | < \infty .
\]
\item The Fourier transform operator \(\Fourier\) is normalized by
\[
\FT{\varphi} (\xi) = (\Fourier \varphi) (\xi)
\eqd \int_{\R^d} \varphi (x) e^{- 2 \pi ix \cdot \xi} \dd x .
\]
\end{itemize}

\begin{definition}[Cancellation classes] \label{def:psi-classes} The set \(\Psi (\R^d) \subset \Sch (\R^d)\) is the set of Schwartz functions whose Fourier transform is supported away from the origin:
\begin{equation} \label{eq:psi-fourier-support}
\Psi (\R^d) \eqd \left\{ f \in \Sch (\R^d) \st 0 \not\in \spt \FT{f} \right\}.
\end{equation}
For a nonnegative integer \(m\), the set \(\Psi^m (\R^d) \subset \Sch (\R^d)\) is the set of Schwartz functions with vanishing moments up to order \(m\):

\begin{equation} \label{eq:psi-vanishing-moments}
\Psi^m (\R^d) \eqd \left\{ f \in \Sch (\R^d) \st \int_{\R^d} f(x) P(x) \dd x = 0
\quad \text{for every polynomial } P \text{ with } \deg(P) \le m \right\}.
\end{equation}
In particular, \(\Psi(\R^d) \subset \Psi^m(\R^d)\) for every nonnegative integer \(m\).
\end{definition}

\begin{definition}[H{\"o}lder spaces] \label{def:holder-spaces}
Let \(\alpha > 0\) be noninteger and set \(n\eqd \lfloor \alpha \rfloor\).

For \(f \in C^n(\R^d)\), define the homogeneous H{\"o}lder seminorm by
\begin{equation}
\label{eq:holder-homogeneous-seminorm}
\|f\|_{\dot{\mathcal{C}}^{\alpha}(\R^d)} \eqd \max_{|\gamma|=n} \sup_{x \neq x'} \frac{| \partial^\gamma f(x') - \partial^\gamma f(x) |} {|x'-x|^{\{\alpha\}}}.
\end{equation}
This is a seminorm; it does not see polynomials of degree at most \(n\).

For \(t_0>0\), define the nonhomogeneous H{\"o}lder norm at top scale \(t_0\) by
\begin{equation} \label{eq:holder-nonhomogeneous-top-scale-norm}
\|f\|_{\mathcal{C}^{\alpha}_{t_0}(\R^d)} \eqd \sum_{|\gamma|\le n} t_0^{|\gamma|} \|\partial^\gamma f\|_{L^\infty} + t_0^\alpha \max_{|\gamma|=n} \sup_{\substack{x\neq x' \\ |x'-x|<t_0 }} \frac{| \partial^\gamma f(x')-\partial^\gamma f(x)|} {|x'-x|^{\{\alpha\}}}.
\end{equation}
When \(t_0=1\) this is the usual nonhomogeneous H{\"o}lder norm used in the unit-scale formulation.
\end{definition}
\begin{remark}[Homogeneity of the nonhomogeneous class] Although the nonhomogeneous H{\"o}lder norm is not scale-invariant for a fixed $t_0$, the spaces are homogeneous as a class.
Specifically, spatial scaling simply shifts the top scale parameter.

If we define the rescaled function $f_\lambda(x) = f(x/λ)$ for $\lambda > 0$, applying the chain rule $\partial^\gamma f_\lambda(x) = \lambda^{-|\gamma|} (\partial^\gamma f)( x/λ)$ yields:
\[
\|f_\lambda\|_{\mathcal{C}^{\alpha}_{λ t_0}(\R^d)} = \|f\|_{\mathcal{C}^{\alpha}_{t_0}(\R^d)}.
\]
Thus, while any individual norm depends strictly on a chosen scale, the parameter family $\{\mathcal{C}^{\alpha}_{t_0}\}_{t_0>0}$ behaves covariantly under dilations, preserving homogeneity at the level of the class.
\end{remark}

\begin{definition}[Wavelet Operators] \label{def:wavelet-operators} For $x \in \R^d$ and $t > 0$, let the translation and dilation operators acting on a fixed test function $\varphi \in \Sch(\R^d)$ be defined respectively by
\[
\Tr_x\varphi(y) \eqd \varphi(y-x) \quad \text{and} \quad \Dil_t\varphi(y) \eqd t^{-d}\varphi(y/t).
\]
\begin{enumerate}[(i)] \item \textbf{Wavelet Coefficient Operator:} For a function $f$ of at most polynomial growth, define
\begin{equation} \label{eq:wavelet-coefficient-fixed-test}
\Em_{\varphi}[f](x,t) \eqd \int_{\R^d} f(y)\Tr_x\Dil_t\varphi(y)\dd y.
\end{equation}

\item \textbf{Continuous Reconstruction Operator:} Let $F:\R^{d+1}_+\to\C$.
Whenever the integral is absolutely convergent, define

\begin{equation} \label{eq:reconstruction-operator}
(\dot{\mc{R}}_{\varphi}[F])(x) \eqd \int_0^\infty \int_{\R^d} F(y,t)\Tr_y\Dil_t\varphi(x) \dd y\frac{\dd t}{t}.
\end{equation}

\item \textbf{Top-Scale Synthesis Operator:} For $t_0>0$ and a function $a:\R^d\to\C$, define

\begin{equation} \label{eq:top-scale-synthesis-operator}
(\mc{R}_{\varphi,t_0}[a])(x) \eqd \int_{\R^d} a(y)\Tr_y\Dil_{t_0}\varphi(x)\dd y.
\end{equation}
\end{enumerate}
\end{definition}

\begin{remark} \label{rem:variable-test-functions} One can more generally consider these operators when the test functions are allowed to depend on the space and scale parameters, provided they remain uniformly bounded in the Schwartz topology.
Specifically:
\begin{itemize}
\item In the continuous reconstruction operator \eqref{eq:reconstruction-operator}, the fixed function $\varphi$ can be replaced by a bounded family of test functions $\varphi: \R^{d+1}_+ \to \Sch(\R^d)$, modifying the integrand to use $\varphi_{(y,t)}$.
\item In the top-scale synthesis operator \eqref{eq:top-scale-synthesis-operator}, $\varphi$ can be replaced by a bounded spatial family $\varphi: \R^d \to \Sch(\R^d)$, modifying the integrand to use $\varphi_y$.
\end{itemize}
The standard definitions presented in Definition~\ref{def:wavelet-operators} are recovered by identifying a fixed Schwartz function $\varphi$ with the respective constant families.
\end{remark}